% =============================================================
% preamble.tex
% パッケージ読み込みと日本語組版の基本設定（LuaLaTeX + luatexja）
% =============================================================

% ---- 日本語組版 ----
\usepackage{luatexja}
\usepackage{luatexja-fontspec}
% 既定フォントはシステム同梱の日本語フォントに任せる（luatexja の標準）
% 環境差を避けるため明示的なフォント指定はしない

% ---- レイアウト ----
\usepackage[a5paper,top=20mm,bottom=22mm,left=18mm,right=18mm,headheight=28pt,headsep=6mm,footskip=12mm]{geometry}
\usepackage{fancyhdr}
\usepackage{enumitem}
\usepackage{parskip}

% ---- 図・表 ----
\usepackage{graphicx}
\usepackage{adjustbox}
\usepackage{tikz}
\usetikzlibrary{positioning,arrows.meta,shapes.geometric,fit,backgrounds,calc,shadows.blur}
\usepackage{array}
\usepackage{booktabs}
\usepackage{longtable}
\usepackage{tabularx}

% ---- 囲み・色 ----
\usepackage{xcolor}
\usepackage[most]{tcolorbox}
\tcbuselibrary{skins,breakable}

% ---- リンク・QR ----
\usepackage{hyperref}
\usepackage{qrcode}

% ---- その他 ----
\usepackage{etoolbox}
\usepackage{xparse}

% ---- テーマ（色・TikZ スタイル）をここで読み込む ----
% xcolor と tikz がロードされた後でなければ \definecolor と \tikzset が使えない
% =============================================================
% theme.tex
% 配色・装飾テーマ（落ち着いた学習書デザイン）
% =============================================================

% メインカラー（濃紺系：落ち着き）
\definecolor{themeMain}{HTML}{1F3A5F}
% サブ（青灰）
\definecolor{themeSub}{HTML}{3C5A80}
% アクセント（シアン寄り：リンクや強調）
\definecolor{themeAccent}{HTML}{0E7C86}
% 警告・注意（赤橙を少しくすませる）
\definecolor{themeAlert}{HTML}{B24B2C}
% 情報補足（やや暖かいベージュ）
\definecolor{themeInfoBg}{HTML}{FBF6EC}
% 要点（淡いミント）
\definecolor{themeKeyBg}{HTML}{EAF4F2}
% つまずき（淡い桃）
\definecolor{themePitBg}{HTML}{FBEDE7}
% 比較（淡い藍）
\definecolor{themeCmpBg}{HTML}{ECF0F7}
% 罫線
\definecolor{themeRule}{HTML}{C8D0DA}
% まとめ（淡い黄）
\definecolor{themeSumBg}{HTML}{FFF8DA}

% TikZ のノード既定
\tikzset{
  node distance=10mm and 12mm,
  every node/.style={font=\small},
  boxnode/.style={
    draw=themeMain, thick, fill=white, rounded corners=3pt,
    inner sep=5pt, align=center, minimum height=8mm
  },
  keynode/.style={
    draw=themeAccent, thick, fill=themeKeyBg, rounded corners=3pt,
    inner sep=5pt, align=center, minimum height=8mm
  },
  subnode/.style={
    draw=themeSub, thin, fill=white, rounded corners=2pt,
    inner sep=4pt, align=center, font=\footnotesize
  },
  groupbg/.style={
    draw=themeRule, dashed, fill=themeCmpBg, rounded corners=4pt,
    inner sep=6pt
  },
  arr/.style={-{Latex[length=2mm]}, themeSub, thick},
  arr2/.style={{Latex[length=2mm]}-{Latex[length=2mm]}, themeSub, thick},
}


% ---- ハイパーリンク色 ----
\hypersetup{
  colorlinks=true,
  linkcolor=themeAccent,
  urlcolor=themeAccent,
  citecolor=themeAccent,
  bookmarksnumbered=true,
  pdftitle={図でつながる ITパスポート},
  pdfauthor={図でつながる ITパスポート 編集部}
}

% ---- 見出しスタイル ----
% ltjsbook のクラス既定見出しを利用する（titlesec は ltjsbook と相性が悪いので使わない）。

% ---- ページヘッダ・フッタ ----
% oneside 運用のため、E/O を分けず左右のみで指定する
\pagestyle{fancy}
\fancyhf{}
\fancyhead[R]{\small\nouppercase\leftmark}
\fancyhead[L]{\small 図でつながる ITパスポート}
\fancyfoot[C]{\small\thepage}
\renewcommand{\headrulewidth}{0.4pt}
\renewcommand{\footrulewidth}{0pt}
\renewcommand{\headrule}{\hbox to\headwidth{\color{themeRule}\leaders\hrule height \headrulewidth\hfill}}

% 章タイトルが \leftmark に流れすぎないよう短縮
\renewcommand{\chaptermark}[1]{\markboth{第\thechapter 章\quad #1}{}}

% ---- リスト整形 ----
\setlist[itemize]{leftmargin=1.4em,itemsep=2pt,topsep=4pt}
\setlist[enumerate]{leftmargin=1.6em,itemsep=2pt,topsep=4pt}

% ---- 段落 ----
\setlength{\parindent}{1em}
\setlength{\parskip}{0.4em plus 0.1em minus 0.1em}
